\chapter{Introducción}
\label{chap:introducion}

\lettrine{E}{l} pádel es el segundo deporte más practicado en España, con más de
cinco millones de jugadores federados y miles de clubes repartidos por todo el
territorio nacional. Este crecimiento exponencial ha generado una demanda
creciente de herramientas digitales que permitan a los centros deportivos gestionar
su actividad de forma eficiente, desde la reserva de pistas hasta la organización
de torneos. Sin embargo, muchos clubes siguen dependiendo de sistemas manuales,
hojas de cálculo o aplicaciones genéricas que no están adaptadas a las
particularidades del negocio del pádel.

PadelCenter nace como respuesta a esta necesidad: una aplicación web específicamente
diseñada para la gestión integral de centros de pádel, que abarca la administración
de pistas, la creación y gestión de reservas, y la organización de torneos con
emparejamientos automáticos.

\section{Contexto y motivación}
\label{sec:contexto}

La transformación digital del sector deportivo es una tendencia global impulsada
por la pandemia de COVID-19 y por el auge de los servicios \textit{on-demand}.
Los usuarios esperan poder reservar una pista desde su teléfono móvil en cuestión
de segundos, con confirmación inmediata y sin necesidad de llamar por teléfono
al club.

Desde el punto de vista académico, este proyecto representa una oportunidad de
aplicar en un contexto real las competencias adquiridas durante el Grado en
Ingeniería Informática, en particular las relacionadas con la ingeniería del
software: diseño de sistemas, arquitecturas limpias, desarrollo \textit{contract-first},
pruebas automatizadas y trabajo con tecnologías modernas tanto en backend como
en frontend.

La elección del dominio del pádel responde a una motivación personal del autor,
combinada con la observación de que los sistemas existentes en el mercado presentan
importantes carencias en cuanto a personalización, extensibilidad y calidad del
software subyacente.

\section{Objetivos}
\label{sec:objetivos}

Los objetivos del presente trabajo se dividen en funcionales y técnicos:

\subsection{Objetivos funcionales}

\begin{itemize}
  \item Implementar un sistema de autenticación seguro con roles diferenciados
        (jugador y administrador de centro).
  \item Permitir la gestión de centros deportivos y sus pistas (alta, baja,
        modificación, consulta).
  \item Desarrollar un módulo de reservas con control de conflictos temporales
        y visibilidad según el rol del usuario.
  \item Implementar un módulo de torneos con inscripción de participantes y
        generación automática de emparejamientos.
\end{itemize}

\subsection{Objetivos técnicos}

\begin{itemize}
  \item Diseñar el sistema siguiendo la arquitectura hexagonal, garantizando
        la independencia del dominio de negocio respecto a los detalles de
        infraestructura.
  \item Adoptar un enfoque \textit{contract-first} basado en OpenAPI~3, de modo
        que la especificación de la API sea la fuente de verdad del contrato
        entre frontend y backend.
  \item Garantizar una cobertura de pruebas adecuada mediante tests unitarios,
        de integración y de arquitectura.
  \item Implementar un entorno de desarrollo reproducible mediante Docker Compose.
\end{itemize}

\section{Estructura de la memoria}
\label{sec:estructura}

El presente documento se organiza en los siguientes capítulos:

\begin{description}
  \item[Capítulo~\ref{chap:estadoarte} — Estado del arte:] Se analizan las
    soluciones existentes en el mercado para la gestión de reservas deportivas,
    se evalúan las tecnologías candidatas y se justifica la elección de la pila
    tecnológica adoptada.

  \item[Capítulo~\ref{chap:analisis} — Análisis:] Se presentan los actores del
    sistema, los requisitos funcionales y no funcionales, y los casos de uso
    principales.

  \item[Capítulo~\ref{chap:diseno} — Diseño:] Se describe la arquitectura del
    sistema, el modelo de dominio, el diseño de la API y el esquema de base
    de datos.

  \item[Capítulo~\ref{chap:implementacion} — Implementación:] Se detallan los
    aspectos más relevantes de la implementación del backend y el frontend,
    así como la integración entre ambos.

  \item[Capítulo~\ref{chap:pruebas} — Pruebas:] Se explica la estrategia de
    pruebas adoptada y se presentan los resultados obtenidos.

  \item[Capítulo~\ref{chap:conclusions} — Conclusiones:] Se evalúan los
    resultados obtenidos, se relaciona el trabajo con las competencias de la
    titulación y se proponen líneas de trabajo futuro.
\end{description}
